\documentclass[12pt,a4paper]{article}

\usepackage[utf8]{inputenc}
\usepackage[T1]{fontenc}
\usepackage[french]{babel}
\usepackage{amsmath, amssymb, amsthm}
\usepackage{geometry}
\usepackage{booktabs}
\usepackage{array}
\usepackage{xcolor}
\usepackage{hyperref}

\geometry{margin=2.5cm}

% ── Théorèmes ────────────────────────────────────────────────
\newtheorem{theoreme}{Théorème}
\newtheorem{corollaire}{Corollaire}
\newtheorem{proposition}{Proposition}
\newtheorem{lemme}{Lemme}
\theoremstyle{definition}
\newtheorem{definition}{Définition}
\newtheorem{remarque}{Remarque}

% ── Commandes ────────────────────────────────────────────────
\newcommand{\Z}{\mathbb{Z}}
\newcommand{\N}{\mathbb{N}}
\newcommand{\ord}{\mathrm{ord}}

% ─────────────────────────────────────────────────────────────

\title{%
  \textbf{Structure arithmétique de la famille de premiers}\\[6pt]
  $p = 3m(m+1)+1$, \quad $m = 2^a \cdot 3^b - 1$\\[10pt]
  \large Section indépendante — à intégrer dans le Papier~II
}

\author{Hassane BAKKAOUI}
\date{Avril 2026}

\begin{document}

\maketitle
\tableofcontents
\newpage

% ═══════════════════════════════════════════════════════════════
\section{Introduction}
% ═══════════════════════════════════════════════════════════════

Nous étudions la famille paramétrique de nombres premiers définie par

\begin{equation}\label{eq:famille}
  p = 3m(m+1)+1, \qquad m = 2^a \cdot 3^b - 1,
  \quad a, b \in \N^*.
\end{equation}

Cette famille est remarquable pour deux raisons complémentaires.
Premièrement, la factorisation de $p-1$ est connue par construction :
\[
  p - 1 = 3m(m+1) = 2^a \cdot 3^{b+1} \cdot m,
\]
ce qui rend le \emph{théorème de Pocklington} directement applicable
et fournit une \textbf{preuve de primalité certifiée} sans factorisation
préalable de $p-1$.

Deuxièmement, comme nous le montrons dans cette section,
la famille possède une structure arithmétique modulaire précise,
qui peut être exploitée algorithmiquement pour accélérer la recherche.

Tous les résultats de cette section ont été vérifiés empiriquement
sur les \textbf{1~000 premiers nombres premiers} de la famille,
obtenus par énumération exhaustive sur $a, b \in [1, 600]$.

% ═══════════════════════════════════════════════════════════════
\section{Propriétés modulaires fondamentales}
% ═══════════════════════════════════════════════════════════════

\subsection{Congruence modulo 6}

\begin{theoreme}[Congruence universelle]\label{thm:mod6}
  Pour tout $a, b \in \N^*$, si $p = 3m(m+1)+1$ avec $m = 2^a \cdot 3^b - 1$,
  alors
  \[
    p \equiv 1 \pmod{6}.
  \]
\end{theoreme}

\begin{proof}
  Nous vérifions les deux congruences séparément.

  \medskip
  \textbf{Modulo 2.}
  Pour tout entier $m$, le produit $m(m+1)$ est le produit de deux
  entiers consécutifs, donc pair. Ainsi $3m(m+1) \equiv 0 \pmod{2}$,
  et $p = 3m(m+1)+1 \equiv 1 \pmod{2}$.

  \medskip
  \textbf{Modulo 3.}
  On a $3m(m+1) \equiv 0 \pmod{3}$, donc $p \equiv 1 \pmod{3}$.

  \medskip
  Par le théorème chinois des restes, $p \equiv 1 \pmod{6}$.
\end{proof}

\begin{remarque}
  Ce résultat est indépendant des paramètres $a$ et $b$.
  Il constitue une contrainte structurelle absolue :
  tout candidat $p$ de cette famille satisfait $p \equiv 1 \pmod{6}$,
  ce qui est cohérent avec le fait que tout premier $> 3$
  vérifie $p \equiv 1$ ou $5 \pmod{6}$.
  La famille \eqref{eq:famille} n'occupe que la classe $1 \pmod{6}$.
\end{remarque}

\subsection{Classes interdites modulo 7}

Posons $x = 2^a \cdot 3^b \bmod 7$.
La congruence de $p$ modulo 7 s'exprime en fonction de $x$ seul.

\begin{lemme}\label{lem:pmod7}
  $p \equiv 3(x-1)x + 1 \pmod{7}$, où $x \equiv 2^a \cdot 3^b \pmod{7}$.
\end{lemme}

\begin{proof}
  On a $m = x - 1 \pmod{7}$ et $m+1 \equiv x \pmod{7}$, donc
  $3m(m+1) \equiv 3(x-1)x \pmod{7}$ et $p \equiv 3(x-1)x + 1 \pmod{7}$.
\end{proof}

\begin{lemme}\label{lem:xmod7}
  L'ordre de $2$ dans $(\Z/7\Z)^*$ est $3$, et l'ordre de $3$ est $6$.
  Par conséquent :
  \begin{align*}
    2^a \bmod 7 &\text{ ne dépend que de } a \bmod 3,\\
    3^b \bmod 7 &\text{ ne dépend que de } b \bmod 6.
  \end{align*}
\end{lemme}

\begin{proof}
  Par calcul direct : $2^3 = 8 \equiv 1 \pmod{7}$ et
  $3^6 = 729 \equiv 1 \pmod{7}$, avec $\ord_7(3) = 6$
  (vérifiable en calculant $3^1, \ldots, 3^6 \bmod 7$).
\end{proof}

\begin{theoreme}[Classes interdites par 7]\label{thm:mod7}
  $p \equiv 0 \pmod{7}$ si et seulement si
  $(a \bmod 3,\; b \bmod 6)$ appartient à l'ensemble
  \[
    \mathcal{F}_7 = \{(0,2),\,(0,3),\,(1,0),\,(1,1),\,(2,4),\,(2,5)\}.
  \]
  Autrement dit, ces six classes modulo $(3, 6)$ sont \emph{interdites}
  pour la primalité de $p$.
\end{theoreme}

\begin{proof}
  Par le Lemme~\ref{lem:pmod7}, $p \equiv 0 \pmod 7$ équivaut à
  $3(x-1)x + 1 \equiv 0 \pmod{7}$, soit $3x^2 - 3x + 1 \equiv 0 \pmod{7}$.
  Par inspection sur $x \in \{0,\ldots,6\}$, les solutions sont
  $x \equiv 2$ et $x \equiv 6 \pmod{7}$.

  Par le Lemme~\ref{lem:xmod7}, $x = 2^a \cdot 3^b \bmod 7$ ne dépend
  que de $(a \bmod 3, b \bmod 6)$.
  Le tableau suivant liste toutes les 18 classes et identifie
  celles qui donnent $x \equiv 2$ ou $6$ :

  \begin{center}
  \renewcommand{\arraystretch}{1.2}
  \begin{tabular}{ccccc}
    \toprule
    $a \bmod 3$ & $b \bmod 6$ & $2^a \bmod 7$ & $3^b \bmod 7$
      & $x \bmod 7$ \\
    \midrule
    0 & 0 & 1 & 1 & 1 \\
    0 & 1 & 1 & 3 & 3 \\
    \rowcolor{red!15}
    0 & 2 & 1 & 2 & \textbf{2} \\
    \rowcolor{red!15}
    0 & 3 & 1 & 6 & \textbf{6} \\
    0 & 4 & 1 & 4 & 4 \\
    0 & 5 & 1 & 5 & 5 \\
    \midrule
    \rowcolor{red!15}
    1 & 0 & 2 & 1 & \textbf{2} \\
    \rowcolor{red!15}
    1 & 1 & 2 & 3 & \textbf{6} \\
    1 & 2 & 2 & 2 & 4 \\
    1 & 3 & 2 & 6 & 5 \\
    1 & 4 & 2 & 4 & 1 \\
    1 & 5 & 2 & 5 & 3 \\
    \midrule
    2 & 0 & 4 & 1 & 4 \\
    2 & 1 & 4 & 3 & 5 \\
    2 & 2 & 4 & 2 & 1 \\
    2 & 3 & 4 & 6 & 3 \\
    \rowcolor{red!15}
    2 & 4 & 4 & 4 & \textbf{2} \\
    \rowcolor{red!15}
    2 & 5 & 4 & 6 & \textbf{6} \\
    \bottomrule
  \end{tabular}
  \end{center}

  Les lignes en rouge correspondent à $x \equiv 2$ ou $6 \pmod 7$,
  soit exactement $\mathcal{F}_7$.
\end{proof}

\begin{corollaire}[Efficacité du crible modulaire]
  Sur l'ensemble de toutes les paires $(a \bmod 3, b \bmod 6)$,
  la proportion de candidats éliminés par divisibilité par $7$ est
  \[
    \frac{|\mathcal{F}_7|}{3 \times 6} = \frac{6}{18} = \frac{1}{3}.
  \]
  Le seul premier $q = 7$ élimine exactement un tiers des candidats $(a,b)$.
\end{corollaire}

\begin{remarque}
  Ce résultat justifie théoriquement l'efficacité du crible modulaire
  implémenté dans l'algorithme PrimeQuest :
  chaque petit premier $q$ élimine une fraction calculable des paires $(a,b)$
  par des opérations modulaires de coût $O(1)$,
  sans jamais calculer $p$ en grande précision.
\end{remarque}

% ═══════════════════════════════════════════════════════════════
\section{Analyse empirique des 1 000 premiers de la famille}
% ═══════════════════════════════════════════════════════════════

Nous avons calculé les 1~000 premiers nombres premiers de la famille,
en énumérant les paires $(a,b)$ par ordre croissant de $a+b$
(diagonales du plan discret), pour $a, b \in [1, 600]$.
La recherche a duré \textbf{27 secondes} sur un processeur standard.

\subsection{Distribution de $a$ et $b$}

\begin{center}
\renewcommand{\arraystretch}{1.2}
\begin{tabular}{lrrrrr}
  \toprule
  Paramètre & Minimum & Maximum & Moyenne & Médiane & Écart-type \\
  \midrule
  $a$       & 1  & 581 & 167.0 & 132 & 135.6 \\
  $b$       & 1  & 584 & 143.6 & 108 & 131.2 \\
  $a - b$   & $-552$ & $574$ & 23.3 & 18 & 201.1 \\
  $a/b$     & 0.002 & 525 & 6.39 & 1.34 & 24.9 \\
  \bottomrule
\end{tabular}
\end{center}

La médiane du ratio $a/b$ vaut $1{,}34$, ce qui indique que la zone
$a \approx b$ est la plus fertile pour la famille.
Cette observation est à l'origine de la stratégie de recherche
en zigzag depuis le centre $a_{\text{centre}} = (D/2)/\log_{10}(6)$.

\subsection{Parité}

La répartition de la parité de $a$ et $b$ est la suivante :

\begin{center}
\begin{tabular}{lrr}
  \toprule
  Classe & Effectif & Proportion \\
  \midrule
  $a$ pair     & 571 & 57{,}1\% \\
  $a$ impair   & 429 & 42{,}9\% \\
  \midrule
  $b$ pair     & 494 & 49{,}4\% \\
  $b$ impair   & 506 & 50{,}6\% \\
  \bottomrule
\end{tabular}
\end{center}

Les quatre combinaisons de parité $(a \bmod 2,\, b \bmod 2)$ sont
toutes représentées, avec des proportions proches de $25\%$ chacune.
\textbf{Il n'existe aucune restriction de parité} pour la primalité
dans cette famille.

\subsection{Analyse de $a + b$}

La somme $a + b$ est uniformément distribuée modulo 2, 3, 4 et 6
(toutes les classes représentées à parts quasi-égales).
\textbf{La somme $a + b$ ne révèle aucune structure particulière.}

\subsection{Classes interdites et vérification du Théorème~\ref{thm:mod7}}

Le tableau croisé $(a \bmod 6) \times (b \bmod 6)$
sur les 1~000 premiers révèle six entrées exactement nulles :

\begin{center}
\renewcommand{\arraystretch}{1.2}
\begin{tabular}{c|rrrrrr}
  \toprule
  $a \!\bmod 6 \;\backslash\; b \bmod 6$
    & 0 & 1 & 2 & 3 & 4 & 5 \\
  \midrule
  0 & 89 & 45 & \textbf{0} & \textbf{0} & 32 & 73 \\
  1 & \textbf{0} & \textbf{0} & 27 & 29 & 21 & 34 \\
  2 & 43 & 57 & 29 & 33 & \textbf{0} & \textbf{0} \\
  3 & 35 & 62 & \textbf{0} & \textbf{0} & 52 & 34 \\
  4 & \textbf{0} & \textbf{0} & 30 & 37 & 61 & 42 \\
  5 & 29 & 34 & 46 & 26 & \textbf{0} & \textbf{0} \\
  \bottomrule
\end{tabular}
\end{center}

Ces 12 zéros (6 paires distinctes modulo $(3,6)$) correspondent exactement
à $\mathcal{F}_7$ du Théorème~\ref{thm:mod7}, confirmé sur les 1~000 premiers.

% ═══════════════════════════════════════════════════════════════
\section{Algorithme de recherche — zigzag depuis le centre}
% ═══════════════════════════════════════════════════════════════

\subsection{Principe}

Pour chercher un premier de $D$ chiffres, la contrainte
$\log_{10}(p) \approx D$ impose
\[
  2(a \log_{10} 2 + b \log_{10} 3) + \log_{10} 3 \approx D.
\]
La solution la plus équilibrée ($a \approx b$) est atteinte pour
\[
  a_{\text{centre}} = \frac{D/2}{\log_{10}(2) + \log_{10}(3)}
  = \frac{D/2}{\log_{10} 6}.
\]

L'algorithme explore les valeurs de $a$ dans l'ordre
\[
  a_{\text{centre}},\;
  a_{\text{centre}}+1,\;
  a_{\text{centre}}-1,\;
  a_{\text{centre}}+2,\;
  a_{\text{centre}}-2,\;\ldots
\]
Pour chaque $a$, les valeurs de $b$ valides sont celles satisfaisant
la contrainte de taille à $\pm 12$ chiffres près
(au plus 7 valeurs de $b$ par $a$).

\subsection{Résultats}

\begin{center}
\renewcommand{\arraystretch}{1.2}
\begin{tabular}{lrrrr}
  \toprule
  Taille cible & $a_{\text{centre}}$ & Tests MR & Temps total & Premier trouvé \\
  \midrule
  10~000 chiffres & 6~425 & 19 & 141 s & $a=6411$, $b=6431$ \\
  \bottomrule
\end{tabular}
\end{center}

Le premier de 9~998 chiffres a été trouvé à $\delta = 14$ du centre
($a = a_{\text{centre}} - 14$), confirmant que la zone centrale
est la plus fertile.

\subsection{Accélération par le crible modulaire}

Avant chaque test de Miller-Rabin (coûteux), le candidat $(a,b)$
est soumis à un \emph{crible modulaire} :
pour chaque petit premier $q \leq 10^6$,
on calcule $p \bmod q$ sans calculer $p$ en grande précision :
\[
  x \equiv 2^{a \bmod (q-1)} \cdot 3^{b \bmod (q-1)} \pmod{q},
  \qquad
  p \bmod q \equiv 3(x-1)x + 1 \pmod{q}.
\]
Cette opération coûte quelques microsecondes, contre plusieurs dizaines
de secondes pour un test de Miller-Rabin sur un nombre de 20~000 chiffres.
Le taux d'élimination observé est de l'ordre de \textbf{97--98\%} des candidats.

% ═══════════════════════════════════════════════════════════════
\section{Preuve de primalité par le théorème de Pocklington}
% ═══════════════════════════════════════════════════════════════

\begin{theoreme}[Pocklington]\label{thm:pocklington}
  Soit $p - 1 = F \cdot R$ avec $F > \sqrt{p}$.
  Si pour chaque facteur premier $q \mid F$
  il existe un entier $w$ tel que
  $w^{p-1} \equiv 1 \pmod{p}$ et
  $\gcd(w^{(p-1)/q} - 1,\, p) = 1$,
  alors $p$ est premier.
\end{theoreme}

Pour la famille \eqref{eq:famille}, on a
$p - 1 = 2^a \cdot 3^{b+1} \cdot m$,
et $F = 2^a \cdot 3^{b+1}$ vérifie $F \approx \sqrt{p}$
(puisque $m \approx F$).
Il suffit donc de trouver des témoins pour $q = 2$ et $q = 3$ uniquement.

\textbf{Exemple (premier de 9~998 chiffres) :}
\begin{align*}
  a &= 6411,\quad b = 6431,\\
  F &= 2^{6411} \cdot 3^{6432} \quad (4999 \text{ chiffres}) > \sqrt{p},\\
  q &= 2 \;\to\; w = 7, \qquad q = 3 \;\to\; w = 5.
\end{align*}

% ═══════════════════════════════════════════════════════════════
\section{Synthèse et perspectives}
% ═══════════════════════════════════════════════════════════════

\subsection{Résultats établis dans cette section}

\begin{enumerate}
  \item \textbf{Théorème~\ref{thm:mod6} :}
    $p \equiv 1 \pmod{6}$ pour tout $(a,b)$ —
    propriété structurelle absolue, indépendante des paramètres.

  \item \textbf{Théorème~\ref{thm:mod7} :}
    Les paires $(a \bmod 3, b \bmod 6) \in \mathcal{F}_7$
    donnent $7 \mid p$ — résultat déterministe,
    vérifié sur les 1~000 premiers de la famille.

  \item \textbf{Corollaire :}
    Le premier $q = 7$ élimine exactement $1/3$ des candidats
    par un calcul modulaire de coût $O(1)$.

  \item \textbf{Aucune restriction de parité :}
    Les quatre classes de parité $(a \bmod 2, b \bmod 2)$
    sont toutes représentées.

  \item \textbf{Distribution centrée :}
    La médiane de $a/b$ est $1{,}34$ — la zone $a \approx b$
    est la plus fertile, justifiant la stratégie de zigzag.
\end{enumerate}

\subsection{Questions ouvertes}

Le Théorème~\ref{thm:mod7} s'étend naturellement à tout premier $q$.
Pour chaque $q$, on peut calculer l'ensemble des classes interdites
$\mathcal{F}_q$ et la fraction de candidats éliminés
$|\mathcal{F}_q| / (\ord_q(2) \cdot \ord_q(3))$.
La somme de ces fractions quantifie précisément l'efficacité
théorique du crible modulaire.

Une question ouverte est de déterminer si la distribution des premiers
dans cette famille satisfait une loi de type \emph{Cramér},
c'est-à-dire si le nombre de premiers avec $a + b \leq N$
est asymptotiquement $c \cdot N^2 / \log N$ pour une constante $c > 0$.

% ─────────────────────────────────────────────────────────────
\end{document}
